\section{Introduction}
Steganography aims to hide not only the content of communication but also the fact that communication is happening. This goal is especially hard in text because language is discrete, structured, and statistically constrained. Even small distributional distortions may leave detectable traces, while small decoding mismatches may break message recovery. Classical and provable-security work established the conceptual basis for this problem~\cite{DBLP:conf/crypto/Simmons83,DBLP:journals/jsac/AndersonP98,DBLP:conf/crypto/HopperLA02}. Later neural and LLM-based methods showed that arithmetic or interval coding over language-model token distributions can embed payload bits into fluent text~\cite{DBLP:conf/acl/FangJA17,ziegler-etal-2019-neural,DBLP:conf/acl/ZhangYYH21,DBLP:conf/naacl/LinLZL24}.

In this paper, we study a practical communication setting with three roles: Alice, Bob, and Censor. Alice and Bob share a prompt and a password in advance. Alice feeds the prompt, password, and secret message into Ghostext, and Ghostext outputs a cover text. The Censor is passive and only inspects the transmitted text; if the text looks natural, the message passes. Bob then uses the same prompt and password to recover the original message. This setting is intentionally narrow, but it captures a realistic first step for security-oriented covert communication engineering.

Our primary concern is not generic generation quality but synchronized, auditable recovery behavior. In local LLM pipelines, decode correctness can silently fail when prompt templates drift, tokenizer behavior changes, seeds differ, or protocol-sensitive runtime parameters diverge. For covert messaging, silent misdecode is dangerous. We therefore prioritize deterministic recoverability under matched conditions and fail-closed behavior under mismatch.

Ghostext follows an encrypt-before-hide strategy and then embeds with finite-interval coding over next-token distributions. The intuition is twofold. First, modern LLM next-token distributions approximate natural language generation behavior. Second, embedding should preserve that behavior as much as possible while still carrying payload bits. Encryption makes the payload computationally indistinguishable from random bits, and arithmetic-coding-style mapping consumes available distributional entropy for token selection. Our use of ``almost perfect'' is a scoped engineering claim about this design intent under explicit assumptions; it is not a universal undetectability claim.

This work makes three contributions. We provide a threat-model-explicit protocol and implementation for local LLM text steganography with deterministic decode under matched settings. We introduce a synchronization discipline that combines packet cryptography, candidate policy constraints, retokenization-stability checks, deterministic quantization, and configuration fingerprints. Finally, we provide a rerunnable evaluation focused on security-relevant questions: matched recoverability, fail-closed mismatch behavior, and practical overhead.
